% Part: first-order-logic
% Chapter: axiomatic-deduction
% Section: soundness

\documentclass[../../../include/open-logic-section]{subfiles}

\begin{document}

\iftag{FOL}
      {\olfileid[fr]{fol}{axd}{sou}}
      {\olfileid[fr]{pl}{axd}{sou}}

\olsection{Correction}

\begin{explain}
Un système de !!{derivation}, tel que la déduction axiomatique, est
\emph{correct} s'il ne permet pas de !!{derive} des résultats qui ne
valent pas réellement. La correction constitue ainsi une garantie de
fiabilité pour les systèmes de !!{derivation}. Selon la propriété de
théorie de la démonstration considérée, nous voulons notamment savoir
que :
\begin{enumerate}
\item toute !!{formula}~$!A$ !!{derivable} est valide;
\item si $!A$ est !!{derivable} à partir d'autres !!{formula}s~$\Gamma$,
  elle en est également une conséquence;
\item si un ensemble de !!{formula}s~$\Gamma$ est incohérent, il est
  insatisfaisable.
\end{enumerate}
Ces propriétés sont essentielles pour un système de !!{derivation}.
Si l'une d'elles était en défaut, le système serait défectueux : il
permettrait de !!{derive} trop de résultats. Il est donc primordial
d'établir la correction d'un système de !!{derivation}.
\end{explain}

\begin{prop}
Si $!A$ est un axiome, alors
\iftag{FOL}
  {$\Sat{M}{!A}[s]$ pour toute !!{structure}~$\Struct{M}$ et toute
    assignation~$s$.}
  {$\pSat{v}{!A}$ pour toute !!{valuation}~$\pAssign{v}$.}
\end{prop}

\begin{proof}
\iftag{FOL}{Il faut vérifier que tous les axiomes sont valides.
  Traitons, par exemple, le cas de l'\olref[qua]{ax:q1}. Supposons que
  $t$ soit !!{free for} $x$ dans~$!A$ et que
  $\Sat{M}{\lforall[x][!A]}[s]$. Par définition de la satisfaction,
  $\Sat{M}{!A}[s']$ pour toute $\varAssign{s'}{s}{x}$, et en
  particulier lorsque $s'(x) = \Value{t}{M}[s]$. D'après
  la \olref[syn][ext]{prop:ext-formulas},
  $\Sat{M}{\Subst{!A}{t}{x}}[s]$. Ainsi,
  $\Sat{M}{(\lforall[x][!A] \lif \Subst{!A}{t}{x})}[s]$.
  L'\olref[qua]{ax:q2} se vérifie de la même façon, et les
  axiomes propositionnels au moyen de leurs tables de vérité.}
  {Pour chaque axiome, sa table de vérité vérifie qu'il s'agit d'une
  tautologie.}
\end{proof}

\begin{thm}[Correction]
\ollabel{thm:soundness}
Si $\Gamma \Proves !A$, alors $\Gamma \Entails !A$.
\end{thm}

\begin{proof}
Procédons par récurrence sur la longueur de la !!{derivation} de~$!A$
à partir de~$\Gamma$. S'il n'y a aucune étape justifiée par une
inférence, toutes les !!{formula}s de la !!{derivation} sont soit des
instances d'axiomes, soit des éléments de~$\Gamma$. D'après la
proposition précédente, tous les axiomes sont
\iftag{FOL}{valides}{des tautologies}; si $!A$ est un axiome, alors
$\Gamma \Entails !A$. Si $!A \in \Gamma$, alors
$\Gamma \Entails !A$ trivialement.

Si la dernière étape de la !!{derivation} de~$!A$ est justifiée par
modus ponens, il existe dans la !!{derivation} des !!{formula}s $!B$
et $!B \lif !A$. L'hypothèse de récurrence s'applique aux parties de
la !!{derivation} qui se terminent par ces !!{formula}s, puisqu'elles
comportent au moins une étape d'inférence de moins. Nous avons donc
$\Gamma \Entails !B$ et $\Gamma \Entails !B \lif !A$. Il s'ensuit
$\Gamma \Entails !A$ d'après
\iftag{FOL}
  {le \olref[syn][sem]{thm:sem-deduction}}
  {le \olref[pl][syn][sem]{thm:sem-deduction}}.

\iftag{FOL}{Supposons maintenant que la dernière étape soit justifiée
par~\QR. Cette étape est de la forme
$!C \lif \lforall[x][!B(x)]$, et une étape antérieure est
$!C \lif !B(c)$, où $c$ ne figure ni dans~$\Gamma$, ni dans~$!C$,
ni dans~$\lforall[x][!B(x)]$. Par hypothèse de récurrence,
$\Gamma \Entails !C \lif !B(c)$. D'après
le \olref[syn][sem]{thm:sem-deduction},
$\Gamma \cup \{!C\} \Entails !B(c)$.

Considérons une structure~$\Struct{M}$ telle que
$\Sat{M}{\Gamma \cup \{!C\}}$. Il faut montrer
$\Sat{M}{\lforall[x][!B(x)]}$. Puisque
$\lforall[x][!B(x)]$ est !!a{sentence}, il suffit de montrer que
$\Sat{M}{!B(x)}[s]$ pour toute assignation de variables~$s$, d'après
la \olref[syn][ass]{prop:sat-quant}. Comme
$\Gamma \cup \{!C\}$ ne contient que des énoncés,
$\Sat{M}{!D}[s]$ pour tout $!D \in \Gamma$, par la
\olref[syn][sat]{defn:satisfaction}. Soit $\Struct{M'}$ la structure
identique à~$\Struct{M}$, sauf que
$\Assign{c}{M'} = s(x)$. Puisque $c$ ne figure ni dans~$\Gamma$ ni
dans~$!C$, le \olref[syn][ext]{cor:extensionality-sent} donne
$\Sat{M'}{\Gamma \cup \{!C\}}$. Comme
$\Gamma \cup \{!C\} \Entails !B(c)$, nous avons
$\Sat{M'}{!B(c)}$. Puisque $!B(c)$ est !!a{sentence}, la
\olref[syn][ass]{prop:sentence-sat-true} donne
$\Sat{M'}{!B(c)}[s]$. D'après
la \olref[syn][ext]{prop:ext-formulas},
$\Sat{M'}{!B(x)}[s]$ si et seulement si
$\Sat{M'}{!B(c)}[s]$; rappelons que $!B(c)$ abrège
$\Subst{!B(x)}{c}{x}$. Ainsi, $\Sat{M'}{!B(x)}[s]$. Comme $c$ ne
figure pas dans~$!B(x)$, la \olref[syn][ext]{prop:extensionality}
donne $\Sat{M}{!B(x)}[s]$. L'assignation~$s$ étant arbitraire,
$\Sat{M}{\lforall[x][!B(x)]}$. Par conséquent,
$\Gamma \cup \{!C\} \Entails \lforall[x][!B(x)]$, puis
$\Gamma \Entails !C \lif \lforall[x][!B(x)]$ d'après
le \olref[syn][sem]{thm:sem-deduction}.

Nous laissons en exercice le cas où $!A$ est justifiée par~\QR, mais
est de la forme $\lexists[x][!B(x)] \lif !C$.}{}
\end{proof}

\tagprob{FOL}
\begin{prob}
Achever la preuve du \olref[fol][axd][sou]{thm:soundness}.
\end{prob}
\tagendprob

\begin{cor}
\ollabel{cor:weak-soundness}
Si $\Proves !A$, alors $!A$ est
\iftag{FOL}{valide}{une tautologie}.
\end{cor}

\begin{cor}
\ollabel{cor:consistency-soundness}
Si $\Gamma$ est satisfaisable, alors il est cohérent.
\end{cor}

\begin{proof}
Démontrons la contraposée. Supposons que $\Gamma$ ne soit pas
cohérent. Alors $\Gamma \Proves \lfalse$, c'est-à-dire qu'il existe
une !!{derivation} de~$\lfalse$ à partir de~$\Gamma$. D'après
le \olref{thm:soundness}, toute
\iftag{FOL}{!!{structure}~$\Struct{M}$}{!!{valuation}~$\pAssign{v}$}
qui satisfait~$\Gamma$ doit satisfaire~$\lfalse$. Or
\iftag{FOL}{$\Sat/{M}{\lfalse}$ pour toute
  !!{structure}~$\Struct{M}$}{$\pSat/{v}{\lfalse}$ pour toute
  !!{valuation}~$\pAssign{v}$}. Aucune
\iftag{FOL}{structure~$\Struct{M}$}{valuation~$\pAssign{v}$} ne peut
donc satisfaire~$\Gamma$; ainsi, $\Gamma$ n'est pas satisfaisable.
\end{proof}

\iftag{FOL}{\begin{editorial}
Dans le cas de la règle~\QR, la source omet le marqueur de formule
devant deux occurrences de~$B$. Il a été rétabli dans
$\lforall[x][!B(x)]$ et $!B(c)$; aucune formule mathématique n'est
modifiée.
\end{editorial}}{}

\end{document}
